\section{Introduction}
Major depressive disorder (MDD) is a leading cause of global disability, affecting millions of individuals worldwide. A substantial subset of patients, estimated at 20--30\%, do not achieve remission despite multiple adequate trials of pharmacological agents, psychotherapy, and electroconvulsive therapy. These patients are classified as having treatment-resistant depression (TRD), a condition associated with significant morbidity, decreased quality of life, and an elevated risk of suicide \citep{Vance2023}.

To address this clinical gap, deep brain stimulation (DBS) has emerged as a promising neurosurgical intervention. Open-loop DBS systems, which deliver constant electrical stimulation to targeted brain structures such as the subcallosal cingulate (SCC) cortex or the ventral capsule/ventral striatum (VC/VS), have demonstrated clinical benefit in open-label studies. However, large-scale randomized controlled trials (RCTs) have shown mixed results, often failing to demonstrate a statistically significant difference between active and sham stimulation groups in short-term phases \citep{Thorne2024}. This discrepancy has been largely attributed to the heterogeneous nature of MDD and the lack of real-time parameter adaptation to fluctuating clinical and electrophysiological states.

To overcome these limitations, closed-loop DBS systems have been developed. Unlike open-loop systems, closed-loop neuromodulation utilizes chronic intracranial recording to detect electrophysiological biomarkers associated with depressive states. When a specific neural signature is detected, the device automatically delivers tailored stimulation to restore healthy network dynamics \citep{Patel2025}. In recent years, several clinical trials have investigated these adaptive closed-loop systems, utilizing platforms developed by engineering teams at companies like Synapse and Meridian.

To date, a comprehensive review of the clinical efficacy and safety of closed-loop DBS in TRD has not been published. Therefore, this systematic review aims to compile and synthesize available clinical evidence regarding closed-loop DBS for TRD, focusing on response rates, remission rates, targeted brain regions, biomarker paradigms, and adverse event profiles.
